\section{Introduction}
\label{sec:intro}

Non-hydrostatic models on surface- and terrain-following ($\sigma$) layers,
such as SWASH \citep{zijlema2011} and NHWAVE \citep{ma2012}, resolve the
dispersion of short waves with a few vertical layers. Because the pressure is
split into a hydrostatic part, carried by the free surface, and a
non-hydrostatic part, they are much cheaper than free-surface Navier--Stokes
solvers based on the volume of fluid, while they keep the vertical structure of
the flow that depth-integrated Boussinesq models lose. Both codes use
structured grids in the horizontal.

\NHFOAM\ brings this class of model into OpenFOAM \citep{weller1998}. It solves
the incompressible Navier--Stokes equations for water only, on a $\sigma$-mesh
whose top boundary is the free surface $z=\eta(x,y,t)$ and moves with it.
Because it is an OpenFOAM solver, the horizontal mesh can be any
two-dimensional mesh (quadrilaterals, triangles or polygons) extruded into
vertical columns, and the model uses OpenFOAM's parallelism, boundary
conditions, turbulence models, \code{fvOptions} and function objects. A
longer-term aim is to couple it to the volume-of-fluid and two-phase sediment
solvers of OpenFOAM (\code{interFoam}, \code{sedInterFoam}) near structures and
in the surf zone.

The present version has
\begin{itemize}[nosep]
  \item 2DV and 3D meshes, unstructured in the horizontal;
  \item the pressure on the layer interfaces, as in NHWAVE and SWASH, with a
        column (line) preconditioner for the resulting matrix, and the
        cell-centred pressure of OpenFOAM as an alternative;
  \item a semi-implicit free surface with no gravity-wave Courant limit;
  \item wave generation and absorption in relaxation zones;
  \item prescribed seabed motion (bottom uplift, rigid landslides);
  \item depth-induced breaking with the hydrostatic front approximation
        \citep{smit2013}.
\end{itemize}
It has no wetting and drying yet.

This report documents the model and collects its verification and validation.
Section~\ref{sec:model} gives the equations and Section~\ref{sec:numerics} the
numerical method, including the pressure arrangement and the time integration.
Section~\ref{sec:verification} covers the first round of tests: analytical
checks, the dispersion of standing waves, the decay of progressive waves, the
submerged bar, seabed motion and breaking, with NHWAVE and REEF3D::NHFLOW
\citep{bihs2016} run on the same set-ups. Section~\ref{sec:ma2012} follows the
test programme of \citet{ma2012}, the paper that introduced NHWAVE, case by
case and against NHWAVE only. It starts from the simplest benchmarks, uses
each one to find and fix weaknesses of \NHFOAM, and states what the fixes
changed. Section~\ref{sec:conclusions} lists the limitations and the plans.
Appendix~\ref{app:usage} explains how to build the code and set up a case.
